\documentclass[11pt]{article}

% cs250preamble.tex --- shared preamble for CS 250 course notes
%
% This file is \input by main.tex and by each individual module
% file so that each module can still compile standalone. Any
% package or custom-environment change should be made here, not
% in the individual module files.

\usepackage[T1]{fontenc}
\usepackage[margin=1in]{geometry}
\usepackage{amsmath, amssymb, amsthm}
\usepackage{enumitem}
\usepackage{hyperref}
\usepackage{booktabs}
\usepackage{listings}
\usepackage{xcolor}
\usepackage{tikz}
\usetikzlibrary{shapes.geometric, shapes.gates.logic.US, shapes.gates.logic.IEC, arrows.meta, positioning, calc, trees}
\usepackage{bussproofs}
\usepackage{mdframed}

\lstset{
  basicstyle=\ttfamily\small,
  frame=single,
  breaklines=true,
  columns=fullflexible,
  mathescape=true
}

\theoremstyle{definition}
\newtheorem{theorem}{Theorem}[section]
\newtheorem{definition}[theorem]{Definition}
\newtheorem{example}[theorem]{Example}

\hypersetup{
  colorlinks=true,
  linkcolor=blue!60!black,
  urlcolor=blue!60!black
}

% Key-takeaway box: used at the end of major sections and at the end of
% each module to summarise the main points. Expects \item bullets inside.
\newmdenv[
  linewidth=0.8pt,
  linecolor=blue!40!black,
  backgroundcolor=blue!3,
  skipabove=8pt,
  skipbelow=8pt,
  innertopmargin=7pt,
  innerbottommargin=7pt,
  innerleftmargin=9pt,
  innerrightmargin=9pt
]{keytakeawaybox}
\newenvironment{keytakeaway}{%
  \begin{keytakeawaybox}%
  \noindent\textbf{Key takeaways.}%
  \begin{itemize}[leftmargin=1.3em, topsep=2pt, itemsep=1pt, label=\textbullet]%
}{%
  \end{itemize}%
  \end{keytakeawaybox}%
}

% Summary/looking-ahead environment: used once at the end of each module
% to wrap up what the module built and point forward.
\newmdenv[
  linewidth=0.8pt,
  linecolor=gray!60,
  backgroundcolor=gray!5,
  skipabove=8pt,
  skipbelow=8pt,
  innertopmargin=7pt,
  innerbottommargin=7pt,
  innerleftmargin=9pt,
  innerrightmargin=9pt
]{summarybox}

% Sidebar environment: used for tangential asides---historical notes,
% pointers to other areas of math/CS, cross-paradigm remarks---that
% enrich the main narrative without being essential to it. Takes one
% argument: the sidebar's title.
\newmdenv[
  linewidth=0.6pt,
  linecolor=gray!60,
  backgroundcolor=gray!4,
  skipabove=8pt,
  skipbelow=8pt,
  innertopmargin=6pt,
  innerbottommargin=6pt,
  innerleftmargin=9pt,
  innerrightmargin=9pt
]{sidebarbox}
\newenvironment{sidebar}[1]{%
  \begin{sidebarbox}%
  \noindent\textbf{Sidebar: #1.}\par\medskip%
}{%
  \end{sidebarbox}%
}

% Reading-map environment: used at the top of each numbered module to
% indicate Epp coverage, prerequisites, relevant intermezzi, and
% suggested practice problems. Expects four \rmentry{label}{body}
% items inside (or arbitrary paragraphs).
\newmdenv[
  linewidth=0.8pt,
  linecolor=black!55,
  backgroundcolor=yellow!4,
  skipabove=10pt,
  skipbelow=14pt,
  innertopmargin=9pt,
  innerbottommargin=9pt,
  innerleftmargin=11pt,
  innerrightmargin=11pt
]{readingmapbox}
\newcommand{\rmentry}[2]{\par\noindent\textbf{#1}\hspace{0.4em}#2\par}
\newenvironment{readingmap}{%
  \begin{readingmapbox}%
  \noindent{\bfseries\large Reading map}%
  \vspace{2pt}%
  \setlength{\parskip}{4pt plus 1pt}%
}{%
  \end{readingmapbox}%
}


\title{CS 250 --- Module 8: Sequences and Recursion}
\author{}
\date{}

\begin{document}
\maketitle

\begin{center}
\fbox{\begin{minipage}{0.92\textwidth}
\textbf{Reading map.} This module covers \textbf{Epp §5.1--5.3}: sequences, mathematical induction, and applying induction to sum identities. \textbf{Prerequisites:} Module 4 (which introduced natural induction) and the ``More on Natural Induction'' intermezzo (which walked through several inductive proofs in detail). \textbf{Relevant intermezzi:} the natural-induction intermezzo between Modules 4 and 5 is a direct prerequisite; the type-algebra intermezzo (after Module 10) complements this module by showing another side of the ``data definitions $=$ induction principles'' story. \textbf{Practice from Epp:} §5.1 exercises on computing terms of sequences and translating between explicit and recursive definitions; §5.2 exercises on induction proofs of sum identities; §5.3 exercises on induction applied to closed-form computations.
\end{minipage}}
\end{center}

\section{Definition by recursion}
We've been building up our toolkit for a while now---logic, sets, functions, induction---and now we're going to put a lot of it together by talking about \emph{sequences}.

What is a sequence? Informally it's just an ordered list of values. You've seen these before! Arrays in C, lists in Python, those are sequences. The values come in a definite order: there's a first element, a second element, a third, and so on.

\begin{definition}[Sequence]
A \emph{sequence}\index{sequence} is a function from the natural numbers (or some subset of them) to some set. If we write $a_1, a_2, a_3, \ldots$ what we really mean is that there's a function $a : \mathbb{N} \to S$ for some set $S$, and $a_n$ is just the value of that function at input $n$. The subscript is the index, the ``position'' in the list.
\end{definition}

\noindent If you're a programmer, think of $a_n$ as \texttt{a[n]}. It's literally the same idea: indexing into a collection.

How do we actually \emph{define} a sequence? There are two main approaches.

\subsection{Explicit formulas}
The first way is to just give a formula that lets you compute the $n$th term directly from $n$. This is called an \emph{explicit formula}\index{explicit formula} or a \emph{closed-form definition}.

For example, $a_n = 2n + 1$ for $n \geq 0$. That gives us the sequence $1, 3, 5, 7, 9, \ldots$ --- the odd numbers. If you want the 100th term you just plug in $n = 99$ and get $a_{99} = 199$. You don't need to know any other terms.

Let's work through a couple more examples from the textbook.

Consider $b_j = \frac{5-j}{5+j}$ for every integer $j \geq 1$. Let's compute the first four terms:
\begin{align*}
b_1 &= \frac{5-1}{5+1} = \frac{4}{6} = \frac{2}{3} \\[4pt]
b_2 &= \frac{5-2}{5+2} = \frac{3}{7} \\[4pt]
b_3 &= \frac{5-3}{5+3} = \frac{2}{8} = \frac{1}{4} \\[4pt]
b_4 &= \frac{5-4}{5+4} = \frac{1}{9}
\end{align*}

And here's another one: $d_m = 1 + \left(\frac{1}{2}\right)^m$ for every integer $m \geq 0$:
\begin{align*}
d_0 &= 1 + \left(\tfrac{1}{2}\right)^0 = 1 + 1 = 2 \\[4pt]
d_1 &= 1 + \left(\tfrac{1}{2}\right)^1 = 1 + \tfrac{1}{2} = \tfrac{3}{2} \\[4pt]
d_2 &= 1 + \left(\tfrac{1}{2}\right)^2 = 1 + \tfrac{1}{4} = \tfrac{5}{4} \\[4pt]
d_3 &= 1 + \left(\tfrac{1}{2}\right)^3 = 1 + \tfrac{1}{8} = \tfrac{9}{8}
\end{align*}

See how the terms are getting closer and closer to 1? That's a \emph{convergent} sequence, but we'll leave that idea for a different class.

\subsection{Recursive definitions}
The other way to define a sequence is \emph{recursively}\index{recursive definition}: you give the first term (or first few terms) and then a rule that builds each subsequent term from the ones that came before.

Instead of the explicit formula $a_n = 2n+1$ we could equivalently say: $a_0 = 1$ and $a_n = a_{n-1} + 2$ for $n \geq 1$. Same sequence! But a very different way of thinking about it.

The classic example of a recursively defined sequence is the Fibonacci sequence:
\begin{lstlisting}
fib(0) = 0
fib(1) = 1
fib(n) = fib(n-1) + fib(n-2)   for n >= 2
\end{lstlisting}

This should look very familiar! It's \emph{exactly} the same structure as writing a recursive function in Python or C. That's not a coincidence. As we discussed back in Module 4 when we introduced BNF syntax and inductive definitions, recursive definitions are really just another case of defining data inductively. The base case gives you the starting point and the recursive case tells you how to build new data from old data.

In fact, remember how in Module 4 we defined the natural numbers as $\text{Nat} := 0 \mid \text{succ Nat}$? A recursively defined sequence is doing the same thing: defining $a_n$ by cases on whether $n$ is the base case or a successor.

\subsection{Summation and product notation}
There's an important piece of notation we need before we go further: \emph{summation notation}\index{summation notation}. You've probably seen the big sigma before:
\[
\sum_{i=1}^{n} a_i = a_1 + a_2 + a_3 + \cdots + a_n
\]

This is just a compact way of writing ``add up all the terms $a_i$ where $i$ goes from 1 to $n$.'' The variable $i$ is the \emph{index} of summation, the bottom tells you where to start, the top tells you where to stop.

If you're a programmer, this is literally a for-loop:
\begin{lstlisting}[language=Python]
total = 0
for i in range(1, n+1):
    total += a[i]
\end{lstlisting}

That's what $\sum$ means. It's a for-loop.\footnote{More precisely, the body is pure---no side effects, no mutation beyond the accumulator. If you want to be fancy about it, it's a \emph{fold}.}

For example:
\[
\sum_{i=1}^{4} i^2 = 1^2 + 2^2 + 3^2 + 4^2 = 1 + 4 + 9 + 16 = 30
\]

And you can also nest these things or get creative with the bounds. The expression
\[
\sum_{k=0}^{n} (-1)^k (k^3 - 1)
\]
is just alternating addition and subtraction of the terms $(k^3 - 1)$ because $(-1)^k$ flips sign.

Similarly, there's \emph{product notation} with the big pi:
\[
\prod_{i=1}^{n} a_i = a_1 \cdot a_2 \cdot a_3 \cdots a_n
\]

The most famous example of this is the factorial: $n! = \prod_{i=1}^{n} i$. Again, same deal as summation but you're multiplying instead of adding.

\section{Induction on sequences}\index{induction on sequences}
Here's where the fun starts. We introduced the induction principle on natural numbers in Module 4 and worked through proofs about addition in the Intermezzo on Natural Induction. Here we use the same tool to prove things about sequences and sums.

Let's recall the rule:
\begin{lstlisting}
natural induction
 assumes P(0)
         subproof:
           assumes P(n)
           ------------
           shows P(succ n)
--------------------------
shows forall n : Nat. P(n)
\end{lstlisting}

Same rule as before. Nothing new here! The only thing that's changing is what $P(n)$ is going to \emph{be}. In the Intermezzo on Natural Induction, $P(n)$ was something like ``$0 + n = n$'' or ``addition is associative.'' Now $P(n)$ is going to be equations about sums and sequences.

\subsection{The standard template}
Here's how an induction proof on a sum identity typically goes. You want to prove that some formula holds for all $n \geq \text{base}$. The template is:
\begin{enumerate}
\item \textbf{State $P(n)$}: write down exactly what you're trying to prove. This is usually an equation like ``the sum of the first $n$ whatevers equals some closed-form expression.''
\item \textbf{Base case}: verify that $P(\text{base})$ is true by direct computation.
\item \textbf{Inductive hypothesis}: assume $P(k)$ is true for some arbitrary $k \geq \text{base}$.
\item \textbf{Inductive step}: prove $P(k+1)$ using the assumption that $P(k)$ holds. This is where the actual work happens.
\end{enumerate}

Let's see this in action.

\subsection{Worked example: sum of the first n integers}
We want to prove: for all $n \geq 1$,
\[
1 + 2 + 3 + \cdots + n = \frac{n(n+1)}{2}
\]

or in summation notation, $\displaystyle\sum_{i=1}^{n} i = \frac{n(n+1)}{2}$.

\textbf{Let $P(n)$ be the statement}: $\displaystyle\sum_{i=1}^{n} i = \frac{n(n+1)}{2}$.

\textbf{Base case} ($n = 1$): The left side is just $1$. The right side is $\frac{1 \cdot 2}{2} = 1$. They're equal, so $P(1)$ is true.

\textbf{Inductive hypothesis}: Assume $P(k)$ is true for some $k \geq 1$. That is, assume
\[
\sum_{i=1}^{k} i = \frac{k(k+1)}{2}
\]

\textbf{Inductive step}: We need to show $P(k+1)$, which says $\displaystyle\sum_{i=1}^{k+1} i = \frac{(k+1)(k+2)}{2}$.

Here's the key trick, and it's the same every time: split the sum into the part we already know about and the new term.
\begin{align*}
\sum_{i=1}^{k+1} i &= \left(\sum_{i=1}^{k} i\right) + (k+1) \\
&= \frac{k(k+1)}{2} + (k+1) && \text{by the inductive hypothesis!} \\
&= \frac{k(k+1)}{2} + \frac{2(k+1)}{2} \\
&= \frac{k(k+1) + 2(k+1)}{2} \\
&= \frac{(k+1)(k+2)}{2}
\end{align*}

And that last expression is exactly what $P(k+1)$ claims. We're done!

See the pattern? You peel off the last term of the sum, apply the inductive hypothesis to the rest, and then do algebra to show you get the right answer. That ``peel off the last term'' move is the heart of every induction proof on sums.

\subsection{Setting up: sum of squares}
Here's another one that you'll see on the assignment. For all $n \geq 1$:
\[
1^2 + 2^2 + \cdots + n^2 = \frac{n(n+1)(2n+1)}{6}
\]

So $P(n)$ is the statement $\displaystyle\sum_{i=1}^{n} i^2 = \frac{n(n+1)(2n+1)}{6}$.

The base case is easy: $P(1)$ says $1 = \frac{1 \cdot 2 \cdot 3}{6} = 1$. Check.

For the inductive step, you'd assume $P(k)$ and then look at
\[
\sum_{i=1}^{k+1} i^2 = \left(\sum_{i=1}^{k} i^2\right) + (k+1)^2 = \frac{k(k+1)(2k+1)}{6} + (k+1)^2
\]

and then the game is to do the algebra to show this equals $\frac{(k+1)(k+2)(2k+3)}{6}$. It's the same strategy---peel, substitute, simplify---just with hairier algebra. I'll leave the details to you for the assignment but the key insight is: factor out $(k+1)$ from both terms and then simplify what's left.

\subsection{Induction is recursion}
I want to close with something that I think is really important and ties together a lot of what we've been doing in this course.

Induction and recursion are \emph{the same thing}. Seriously. They're two sides of the same coin.

When you write a recursive function you compute a value by reducing to a smaller case:
\begin{lstlisting}[language=Python]
def sum_to(n):
    if n == 0:
        return 0          # base case
    else:
        return sum_to(n-1) + n  # recursive case
\end{lstlisting}

When you do an induction proof you prove a statement by reducing to a smaller case: you prove $P(k+1)$ by using $P(k)$. The inductive hypothesis is playing \emph{exactly} the role of the recursive call. The base case of the proof is the base case of the recursion. The inductive step is the recursive step.

Another way to say this is: a recursive function computes a \emph{value} for each natural number by building up from the base case. An inductive proof computes a \emph{proof} for each natural number by building up from the base case. The structure is identical. Induction \emph{is} recursion on proofs.

This is why, back in Module 4, we introduced induction right alongside BNF definitions and recursive structure. They're not separate topics that happen to be in the same chapter---they're the same idea wearing different hats. And once you see that connection, both recursive programming and inductive proofs become a lot less mysterious. You just ask yourself: what's the base case, and how do I build the next one from the previous one?

\section{Common mistakes}

\begin{enumerate}
  \item \textbf{Forgetting to state $P(n)$ explicitly.} A good induction proof starts by writing down exactly what the claim $P(n)$ is. Students sometimes jump straight to the base case without telling the reader what they're inducting over, which makes it impossible to check whether the inductive step is proving the right thing.
  \item \textbf{Inductive step that doesn't use the inductive hypothesis.} If your inductive step from $P(k)$ to $P(k+1)$ never actually references $P(k)$, something is off. Either the claim can be proved directly without induction (in which case you should use a direct proof), or you've made an error and you're proving a weaker claim than you think.
  \item \textbf{Wrong base case.} The base case has to match the lower bound of the claim. ``For all $n \geq 5$'' wants a base case at $n = 5$, not $n = 0$ and not $n = 1$. Students default to $n = 1$ and hope for the best, which gives vacuous base cases for claims that start higher up.
  \item \textbf{Confusing the sum with the last term.} ``Peel off the last term'' means splitting $\sum_{i=1}^{k+1} a_i$ into $(\sum_{i=1}^{k} a_i) + a_{k+1}$. Students sometimes write the split as $a_k + a_{k+1}$---just the last two terms, omitting the rest of the sum. That's a very common error that usually shows up as the algebra ``not working out.''
  \item \textbf{Indexing off by one.} Is the sum $\sum_{i=1}^{n}$ (starting at 1) or $\sum_{i=0}^{n}$ (starting at 0)? The two differ by whichever term you're including/excluding at the ends, and the closed form shifts accordingly. Always look at the bounds of summation carefully; an off-by-one there will propagate through the whole proof.
  \item \textbf{Treating ``$\ldots$'' as a proof step.} In writing $1 + 2 + 3 + \cdots + n = \frac{n(n+1)}{2}$, the $\cdots$ is shorthand for ``and I promise the pattern continues,'' but it's not a proof. The induction \emph{is} the proof that the pattern really continues. Students sometimes try to do arithmetic ``inside'' the dots---for example, writing $1 + 2 + \cdots + k + (k+1) = 1 + 2 + \cdots + (k+1)$ and calling that a step---which is really just a tautology and contributes nothing. The meaningful step is always the one that applies the inductive hypothesis.
\end{enumerate}

\section{Summary and looking ahead}

\begin{keytakeaway}
  \item A \emph{sequence} is just a function from $\mathbb{N}$ (or a subset) to some set. You can define a sequence explicitly ($a_n = 2n+1$) or recursively ($a_0 = 1$, $a_n = a_{n-1} + 2$).
  \item \emph{Summation notation} ($\sum$) and \emph{product notation} ($\prod$) are compact ways to write for-loops in mathematics.
  \item An induction proof on a sum identity always follows the same template: state $P(n)$, check the base case, assume $P(k)$, and prove $P(k+1)$ by peeling off the last term and applying the IH.
  \item \emph{Induction is recursion on proofs.} A recursive function computes a value; an inductive proof ``computes'' a proof. The base case of the code is the base case of the proof; the recursive call is the inductive hypothesis. This is not a metaphor.
\end{keytakeaway}

\begin{summarybox}
\textbf{What we built this module.} We're now fluent with induction as a proof tool. The sum-formula exercises are finger exercises; the deeper point is that every recursive definition \emph{is} an induction principle in disguise, and every inductive proof \emph{is} a recursive function with ``proof'' as its return type. That's the unification we've been building toward since Module 4.

\medskip
\textbf{Looking ahead.} Module 9 generalizes ordinary induction to \emph{strong induction}, where the inductive hypothesis gives you $P(j)$ for \emph{all} $j \leq k$, not just $j = k$. This turns out to be the right tool for recurrences that reach back more than one step, divide-and-conquer algorithms, and (eventually in Module 10) structural induction on arbitrary recursive data types. Also in Module 9: the characteristic-equation method for solving linear recurrences in closed form, which is how we get the Fibonacci formula with $\varphi$ and $\sqrt{5}$.
\end{summarybox}

\section{Exercises}

Organized into \textbf{warm-up} (computation with sequences and summation), \textbf{standard} (inductive proofs following the template), \textbf{challenge} (tighter arguments and inequalities), and \textbf{connection} (programming and forward pointers). Solutions are in the appendix.

\subsection*{Warm-up}

\textbf{Exercise 1.}\label{ex:m8:1} Write a recursive definition for the sequence $a_n = 3n$ (i.e., give $a_0$ and a rule for $a_n$ in terms of $a_{n-1}$). Verify your definition by computing the first four terms and checking that they match $0, 3, 6, 9$.

\bigskip
\textbf{Exercise 2.}\label{ex:m8:2} Compute the following sums explicitly:
\begin{enumerate}[label=(\alph*)]
  \item $\displaystyle\sum_{i=1}^{5} (2i + 1)$
  \item $\displaystyle\sum_{k=0}^{4} 2^k$
  \item $\displaystyle\sum_{j=1}^{3} (-1)^{j+1} j^2$
  \item $\displaystyle\prod_{i=1}^{4} i$ (yes, this is $4!$)
\end{enumerate}

\bigskip
\textbf{Exercise 3.}\label{ex:m8:3} Translate between the two forms:
\begin{enumerate}[label=(\alph*)]
  \item The sequence $a_n$ is defined by $a_0 = 2$ and $a_n = a_{n-1} + 4$ for $n \geq 1$. Give an explicit formula for $a_n$.
  \item The sequence $b_n = n^2 + 1$ is defined explicitly for $n \geq 0$. Give a recursive definition.
\end{enumerate}

\bigskip
\textbf{Exercise 4.}\label{ex:m8:4} Compute the first five terms of the Fibonacci sequence (starting from $\text{fib}(0) = 0$) and verify by hand that $\text{fib}(5) = \text{fib}(4) + \text{fib}(3)$.

\subsection*{Standard}

\textbf{Exercise 5.}\label{ex:m8:5} Prove by induction: for all $n \geq 1$, $\displaystyle\sum_{i=1}^{n} (2i - 1) = n^2$. (This says the sum of the first $n$ odd numbers is $n^2$.) Follow the template: state $P(n)$, check the base case, assume $P(k)$, and prove $P(k+1)$ by peeling off the last term.

\bigskip
\textbf{Exercise 6.}\label{ex:m8:6} Prove by induction: for all $n \geq 1$, $\displaystyle\sum_{i=1}^{n} i^2 = \frac{n(n+1)(2n+1)}{6}$. Do the full algebra carefully---this is a drill-heavy exercise with hairier arithmetic than the $n$-th-triangular-number proof.

\bigskip
\textbf{Exercise 7.}\label{ex:m8:7} Prove by induction: for all $n \geq 0$, $\displaystyle\sum_{i=0}^{n} 2^i = 2^{n+1} - 1$. (This is the formula for a geometric series with ratio $2$.)

\bigskip
\textbf{Exercise 8.}\label{ex:m8:8} Prove by induction: for all $n \geq 1$, $\displaystyle\sum_{i=1}^{n} i^3 = \left(\frac{n(n+1)}{2}\right)^2$. This striking identity says that the sum of the first $n$ cubes equals the \emph{square} of the $n$-th triangular number---a surprise that's worth sitting with. (Hint: you already know the sum of the first $n$ integers; use that in your algebra.)

\subsection*{Challenge}

\textbf{Exercise 9.}\label{ex:m8:9} Prove by induction: for all $n \geq 1$, $\displaystyle\sum_{i=1}^{n} \frac{1}{i(i+1)} = \frac{n}{n+1}$. (Hint: in the inductive step, after you apply the IH, look for a common denominator. The algebra is surprisingly clean.)

\bigskip
\textbf{Exercise 10.}\label{ex:m8:10} Prove by induction: for all $n \geq 4$, $2^n < n!$. This is an \emph{inequality} rather than an equality, which means the inductive step needs different algebra than the sum proofs. (Hint: the base case is $n = 4$. In the inductive step, assume $2^k < k!$ for $k \geq 4$ and show $2^{k+1} < (k+1)!$. Notice that $2^{k+1} = 2 \cdot 2^k$ and $(k+1)! = (k+1) \cdot k!$, so the IH gives you a factor of $2$ on the left and a factor of $(k+1) \geq 5$ on the right.)

\bigskip
\textbf{Exercise 11.}\label{ex:m8:11} Define a sequence by $T_1 = 1$ and $T_n = T_{n-1} + n$ for $n \geq 2$. (These are the triangular numbers.) Prove by induction that $T_n = \frac{n(n+1)}{2}$. Compare your proof to the worked example in this module---the only difference should be the starting point and the naming.

\subsection*{Connection}

\textbf{Exercise 12.}\label{ex:m8:12} (Programming.) Write a recursive Python function \texttt{sum\_of\_squares(n)} that computes $\displaystyle\sum_{i=1}^{n} i^2$. Then compare its structure to the induction proof of the sum-of-squares formula (Exercise~6). Which part of your code corresponds to the base case? Which part corresponds to the inductive step? Where does the ``inductive hypothesis'' show up in the code?

\bigskip
\textbf{Exercise 13.}\label{ex:m8:13} (Programming.) Write both a recursive and an iterative Python function that computes the $n$-th Fibonacci number. Test that they agree on $n = 0, 1, 2, \ldots, 10$. Then time both versions on $n = 30$. Which is faster, and by how much? Why? (If you want, also write a memoized version and time that.)

\bigskip
\textbf{Exercise 14.}\label{ex:m8:14} (Intermezzo pointer.) The intermezzo on type algebra (after Module~10) explains that BNF-style recursive data definitions correspond to \emph{initial algebras} of functors, and that induction principles are \emph{unique-morphism} properties of these algebras. As a preview: the natural numbers $\text{Nat} := 0 \mid \text{succ Nat}$ have a single induction principle, and so does each list type, each tree type, etc. Without peeking ahead, speculate on why ``one induction principle per recursive type definition'' is such a clean correspondence. What does it suggest about how to define induction principles for new types you might encounter later?

\end{document}
